\section{Introduction}
\label{sec:introduction}

LLM-based query expansion has become a common approach to improving retrieval
effectiveness \citep{gao-etal-2023-precise,wang-etal-2023-query2doc,zhang-etal-2024-exploring-best}.
For hybrid retrieval, its gains are typically measured at a preset top-$L$
cutoff using metrics such as nDCG and recall
\citep{zhang-etal-2024-exploring-best,liu-zhang-2025-exp4fuse}.
Recent work on fixed top-$L$ fusion shows that $L$ does more than control access
depth: changing the cutoff can also change which candidates and cross-channel
contributions enter the fused ranking \citep{zhang2026exact}.
Consequently, a gain observed at one $L$ reflects both the query expansion and
the chosen cutoff; at another $L$, its magnitude---and even the ordering between
methods---may change. Evidence from a single cutoff therefore does not isolate
the effect of query expansion.

To remove this confound, we evaluate the retrieval effectiveness of each query
construction using complete-list fusion, without truncating either channel
\citep{zhang2026exact}. We then reveal the dense and sparse rankings
progressively and record the per-channel depths at which the same ordered
top-$K$ becomes fixed. This separates two questions: what fused ranking does an
expansion produce, and how much ranked evidence is required to determine it?
The answers need not move together. An expansion may improve retrieval
effectiveness while pushing decisive evidence deeper into one channel. We
therefore ask whether query expansion can improve effectiveness without
increasing the access depth of either channel.

Meeting both goals requires treating dense and sparse retrieval differently.
Prior studies show that expansion gains vary across retrievers, and that dense
and sparse retrieval can favor different forms of generated evidence
\citep{weller-etal-2024-generative,li-etal-2026-dvcqr}. Our NFCorpus analysis
shows what happens when one shared expansion is applied to both channels. When
the same generated passages are used in both channels, retrieval effectiveness
improves and fewer dense ranks are needed to determine the fused result, yet
more sparse ranks are required. The unexpanded sparse query matches only a
small portion of the corpus; the added terms make far more documents eligible
for sparse ranking, so the list no longer exhausts early. The same generated
content therefore supplies useful semantic evidence to the dense retriever
while spreading lexical evidence across a much larger sparse list. A shared
expansion cannot control these two effects separately.

We address this asymmetry with channel-asymmetric query expansion. The LLM
generates several complementary reference passages around the original query.
For dense retrieval, orthogonal residual expansion preserves the original query
direction and adds only the new semantic component supplied by the references.
For sparse retrieval, score-product anchoring multiplies each document's scores
under the unexpanded and expanded queries. Documents matched only by generated
terms receive no anchored score, while the new terms can still reorder
documents supported by the original query. Together, these operations broaden
semantic coverage without broadening sparse support: Dense expands; sparse
anchors.

The evaluation mirrors this argument. A top-$L$ sensitivity study first
examines whether the measured gains and relative ranking of query expansion
methods change with the cutoff. We then evaluate complete-list fusion on seven
public retrieval benchmarks, comparing our method with the unexpanded query,
shared expansion, and representative generative expansion methods in terms of
retrieval effectiveness and per-channel access depth. Controlled experiments
examine the contributions of complementary reference generation, orthogonal
residual expansion, and score-product anchoring, followed by robustness
analyses across generation and retrieval settings. Together, these experiments
test whether channel-asymmetric query expansion can produce better fused
rankings while reducing the access required from both channels.
